\chapter{Exponential Families}
\label{cha:exp}

Exponential families form a broad and important class of statistical models with many desirable
theoretical and computational properties. They provide a unifying framework for a wide range of
commonly used distributions and play a central role in statistical inference.

A key reason for the popularity of exponential family models is their favorable optimization
structure. In many cases, parameter estimation reduces to maximizing a concave log-likelihood
(or equivalently, minimizing a convex objective function), which leads to stable and efficient
numerical algorithms.

Several widely used models arise naturally as members of exponential families. For example,
logistic regression for binary data is based on a specific parametric form that links predictors
to event probabilities through log-odds. This formulation is not only interpretable but also
ensures that estimation involves a convex optimization problem. Similarly, the Poisson
distribution is a standard choice for modeling count data, partly due to its mathematical
convenience and tractable likelihood.

More generally, exponential families provide a common structure underlying these models.
Many results that are traditionally derived on a case-by-case basis can be obtained more
systematically by working within this general framework. Throughout this course, exponential
families will serve as an important source of examples and will be revisited in various contexts.

We now introduce the formal definition of exponential families and discuss their main properties.

\section{Definition}%
\label{sec:Definition}

\begin{definition}[]
    Let \(\cc{P}=\{\mathsf{P}_{\theta} \::\: \theta\in \Theta\}\) be a statistical model on
    \((\cc{X},\cc{A})\) that is dominated by a \(\sigma\)--finite measure \(\nu\).
    The model \(\cc{P}\) is called an \textbf{exponential family} if there exist an integer
    \(k \in \N\) and functions
    \begin{align*}
        B : \Theta \to \R, \quad \quad 
        &\eta : \Theta \to \R^{k},\\
        T : \cc{X}\to\R^{k}, \quad \quad 
        &h:\cc{X}\to[0,\infty),
    \end{align*}
    such that each distribution \(\mathsf{P}_{\theta}\) admits a density
    \(p_{\theta}(x) = \frac{d \mathsf{P}_{\theta}}{d \nu}(x)\) of the form
    \begin{equation}
        \label{eq:exp}
        p_{\theta}(x)
        =
        \exp\!\left\{
            \left\langle \eta(\theta), T(x) \right\rangle
            - B(\theta)
        \right\} h(x),
        \quad x \in \cc{X}
    \end{equation}
    where 
    \[
        \left<\eta(\theta),T(x) \right> = \sum\limits_{i=1}^{k} \eta_{i}(\theta)T_{i}(x). 
    \]
    
\end{definition}

This representation expresses the log--density as a linear function of the statistic \(T(x)\),
with coefficients \(\eta(\theta)\) depending on the parameter. In particular,
\[
\log p_{\theta}(x)
=
\left\langle \eta(\theta), T(x) \right\rangle
- B(\theta)
+ \log h(x),
\]
which simplifies likelihood-based inference by reducing it to linear and convex-analytic
operations.

The map \(\eta(\theta)\) is called the \textbf{natural parameter}, and \(T\) is the
\textbf{sufficient statistic} of the family (by Neyman's factorization criterion).
The function \(B(\theta)\) acts as a normalizing term ensuring that \(p_{\theta}\) integrates
to one, while \(h(x)\) modifies the dominating measure \(\nu\).

The integer \(k\) is referred to as the \textbf{dimension} of the representation and is, in
general, not unique (because the dimension of sufficient statistics of a family aren't). A statistical model is a \(k\)--dimensional exponential family if all its
densities admit a representation of the form \eqref{eq:exp} with the same statistic \(T\).

\begin{ex}[Binomial model]
    Let \(X \sim \text{Binomial}(n,\theta)\) with \(\theta \in (0,1)\).
    The distribution of \(X\) admits a density with respect to the counting measure on
    \(\cc{X}=\{0,1,\dots,n\}\) given by
    \[
        p_{\theta}(x)
        =
        \binom{n}{x}\theta^{x}(1-\theta)^{n-x}.
    \]

    Writing the parameter--dependent part on the log scale yields
    \[
        p_{\theta}(x)
        =
        \binom{n}{x}
        \exp\!\left\{
            x \log \theta + (n-x)\log(1-\theta)
        \right\},
    \]
    which is of exponential family form. This representation corresponds to a
    \(2\)--dimensional exponential family with sufficient statistic
    \(T(x) = (x, n-x)\) and natural parameters
    \(\eta(\theta) = (\log\theta, \log(1-\theta))\).

    The representation is, however, redundant since the components of \(T(x)\) satisfy
    the linear relation \(x + (n-x) = n\). Using this relation, the density can be rewritten as
    \[
        p_{\theta}(x)
        =
        \binom{n}{x}
        \exp\!\left\{
            x \log\frac{\theta}{1-\theta}
            + n \log(1-\theta)
        \right\}.
    \]

    This yields a \(1\)--dimensional representation of the Binomial model, with sufficient
    statistic \(T(x)=x\) and natural parameter
    \[
        \eta(\theta) = \log\frac{\theta}{1-\theta}.
    \]

The quantity \(\theta/(1-\theta)\) is called the \emph{odds} associated with the success
probability \(\theta\). More generally, for an event \(A\) with probability
\(P(A)=\theta\), the odds are defined as the ratio
\[
    \frac{P(A)}{P(A^{c})}
    =
    \frac{\theta}{1-\theta}.
\]
While probabilities take values in \((0,1)\), odds range over \((0,\infty)\).
Taking logarithms yields the \emph{log--odds}
\[
    \log\frac{\theta}{1-\theta},
\]
which form the natural parameter \(\eta(\theta)\) of the Binomial exponential family in its
one--dimensional representation. 
%This explains why log--odds play a central role
%in models for binary data, such as logistic regression.
    This representation highlights the connection between
    the Binomial model and generalized linear models, in particular logistic regression.
\end{ex}

\begin{ex}[Multinomial model as an exponential family]
Consider \(n\) independent repetitions of a random experiment with \(s\) possible outcomes.
Let \(X_i \ge 0\) denote the number of times outcome \(i\) is observed, and define
\[
\bs{X} = (X_1,\dots,X_s), \qquad \sum_{i=1}^s X_i = n.
\]
Assume that the probability of outcome \(i\) in a single trial is \(\theta_i>0\), with
\(\sum\limits_{i=1}^s \theta_i = 1\).
Then
\[
\bs{X} \sim \text{Multinomial}(n,\theta_1,\dots,\theta_s),
\]
with probability mass function
\[
p_{\bs{\theta}}(x_1,\dots,x_s)
= \frac{n!}{\prod_{i=1}^s x_i!}\prod_{i=1}^s \theta_i^{x_i}.
\]

Writing the density in exponential form yields
\[
p_{\bs{\theta}}(x_1,\dots,x_s)
= \exp\!\left\{\sum_{i=1}^s x_i \log(\theta_i)\right\}
\cdot \frac{n!}{\prod_{i=1}^s x_i!},
\]
which shows that the multinomial model is an exponential family with sufficient statistic
\(T(x)=(x_1,\dots,x_s)\).
This representation, however, is redundant, since the counts satisfy the linear constraint
\(\sum_{i=1}^s X_i=n\).

To obtain a minimal representation, write
\[
X_s = n - \sum_{i=1}^{s-1} X_i
\]
and substitute into the exponent to obtain
\[
p_{\bs{\theta}}(x_1,\dots,x_s)
= \exp\!\left\{
\sum_{i=1}^{s-1} x_i \log\!\left(\frac{\theta_i}{\theta_s}\right)
+ n \log(\theta_s)
\right\}
\cdot \frac{n!}{\prod_{i=1}^s x_i!}.
\]

This shows that the multinomial model forms an \((s-1)\)-dimensional exponential family with
sufficient statistic
\[
T(x) = (x_1,\dots,x_{s-1})
\]
and natural parameter vector
\[
\eta(\bs{\theta})
= \left(
\log\!\left(\frac{\theta_1}{\theta_s}\right),
\dots,
\log\!\left(\frac{\theta_{s-1}}{\theta_s}\right)
\right).
\]
The term \(\frac{n!}{\prod_{i=1}^s x_i!}\) plays the role of the base measure \(h(x)\) and does
not depend on the parameter \(\bs{\theta}\).
\end{ex}

\begin{ex}[Gaussian model]
Let \(X \sim \cc{N}(\mu,\sigma^2)\) with unknown parameter
\(\bs{\theta}=(\mu,\sigma^2)\in\R\times(0,\infty)\).
The density on \(\cc{X}=\R\) is
\[
p_{\bs{\theta}}(x)
= \frac{1}{\sqrt{2\pi\sigma^2}}
\exp\!\left\{-\frac{1}{2\sigma^2}(x-\mu)^2\right\}.
\]

Expanding the quadratic term yields
\[
p_{\bs{\theta}}(x)
= \frac{1}{\sqrt{2\pi\sigma^2}}
\exp\!\left\{
-\frac{1}{2\sigma^2}x^2
+ \frac{\mu}{\sigma^2}x
- \frac{\mu^2}{2\sigma^2}
\right\}.
\]
The terms depending on the data \(x\) appear linearly as \(x\) and \(x^2\), which allows the
density to be written in exponential family form. Indeed,
\[
p_{\bs{\theta}}(x)
= \exp\!\left\{
\left\langle
\begin{pmatrix} x \\[0.2em] -\frac{x^2}{2} \end{pmatrix},
\begin{pmatrix} \frac{\mu}{\sigma^2} \\[0.2em] \frac{1}{\sigma^2} \end{pmatrix}
\right\rangle
- B(\bs{\theta})
\right\},
\]
where
\(
B(\bs{\theta}) = \frac{\mu^2}{2\sigma^2} + \frac{1}{2}\log(2\pi\sigma^2)
\)
is the log-normalizing constant.

Hence, the Gaussian model is a two-dimensional exponential family with
\begin{itemize}
    \item sufficient statistic \(T(x) = \bigl(x,\,-\tfrac{x^2}{2}\bigr)\),
    \item natural parameter \(\eta(\bs{\theta}) = \bigl(\tfrac{\mu}{\sigma^2},\,\tfrac{1}{\sigma^2}\bigr)\).
\end{itemize}
In this representation the base measure \(h(x)\) is constant, corresponding to Lebesgue measure.
\end{ex}

\myparagraph{Why is not the Gaussian model one-dimensional?}
Although a single observation \(X\) uniquely determines the statistic
\(T(X) = \bigl(X,-\tfrac{X^2}{2}\bigr)\), the Gaussian model cannot be represented
as a one-dimensional exponential family.
The reason is geometric: the map
\[
x \longmapsto \left(x,-\tfrac{x^2}{2}\right)
\]
traces out a nonlinear curve in \(\R^2\).
Consequently, there is no nontrivial linear functional that recovers both
components of \(T(X)\) from a single scalar statistic while preserving the
required inner-product structure of an exponential family.

This contrasts with the binomial case, where the corresponding embedding
\(
x \mapsto (x,n-x)
\)
lies on an affine one-dimensional subspace, allowing a reduction to a
one-dimensional representation.
For the Gaussian model, the quadratic dependence on the data is essential and
cannot be absorbed into a single sufficient statistic without leaving the
exponential family framework.

\begin{ex}[Multivariate Gaussian model]
Let \(\bs{X} \sim \cc{N}_{k}(\bs{\mu},\Sigma)\), where the parameter
\[
\theta = (\bs{\mu},\Sigma) \in \R^{k} \times PD(k),
\]
and \(PD(k)\) denotes the cone of symmetric positive definite \(k\times k\) matrices.
The density of \(\bs{X}\) with respect to Lebesgue measure on \(\cc{X}=\R^{k}\) is
\[
p_{\theta}(\bs{x})
= \frac{1}{\sqrt{(2\pi)^{k}\det(\Sigma)}}
\exp\!\left\{
-\frac{1}{2}(\bs{x}-\bs{\mu})^{T}\Sigma^{-1}(\bs{x}-\bs{\mu})
\right\}.
\]

\medskip
Expanding the quadratic form and discarding terms that do not depend on \(\bs{x}\) yields
\[
p_{\theta}(\bs{x})
\propto
\exp\!\left\{
-\frac{1}{2}\bs{x}^{T}\Sigma^{-1}\bs{x}
+ \bs{\mu}^{T}\Sigma^{-1}\bs{x}
\right\}.
\]

\medskip
The linear term \(\bs{\mu}^{T}\Sigma^{-1}\bs{x}\) is an inner product in \(\R^{k}\) between
\(\bs{x}\) and the vector \(\Sigma^{-1}\bs{\mu}\).
To express the quadratic term in inner-product form, write
\[
\bs{x}^{T}\Sigma^{-1}\bs{x}
= \sum_{i=1}^{k}\sum_{j=1}^{k} x_i x_j (\Sigma^{-1})_{ij}.
\]
Introducing the outer product \(\bs{x}\bs{x}^{T}\), whose \((i,j)\)-entry equals \(x_i x_j\),
this can be rewritten as
\[
\bs{x}^{T}\Sigma^{-1}\bs{x}
= \sum_{i,j} (\Sigma^{-1})_{ij} (\bs{x}\bs{x}^{T})_{ij}
= \tr\!\left( \Sigma^{-1}\bs{x}\bs{x}^{T} \right).
\]

\medskip
The trace induces an inner product on the vector space of symmetric \(k\times k\) matrices,
defined by
\[
\langle A,B\rangle := \tr(AB),
\qquad A,B \in \mathrm{Sym}(k).
\]
Since both \(\Sigma^{-1}\) and \(\bs{x}\bs{x}^{T}\) are symmetric, the quadratic term is an
inner product in this space.

\medskip
Combining both parts, the density can be written in exponential-family form as
\[
p_{\theta}(\bs{x})
\propto
\exp\!\left\{
\left\langle \Sigma^{-1}\bs{\mu},\, \bs{x} \right\rangle
+
\left\langle \Sigma^{-1},\, -\tfrac12 \bs{x}\bs{x}^{T} \right\rangle
\right\}.
\]

\medskip
Thus, the multivariate Gaussian model is an exponential family with
\begin{itemize}
    \item sufficient statistic
    \(
    T(\bs{x}) = \left(\bs{x}, -\tfrac12 \bs{x}\bs{x}^{T}\right),
    \)
    \item natural parameter
    \(
    \eta(\theta) = \left(\Sigma^{-1}\bs{\mu},\, \Sigma^{-1}\right).
    \)
\end{itemize}

\medskip
The dimension of this exponential family equals the dimension of the vector space in which
\(T(\bs{x})\) takes values.
The first component \(\bs{x}\) lies in \(\R^{k}\), contributing \(k\) dimensions.
The second component lies in the space of symmetric \(k\times k\) matrices, whose dimension
equals the number of free entries:
\(k\) diagonal elements and \(\frac{k(k-1)}{2}\) off-diagonal elements.
Hence,
\[
\dim \mathrm{Sym}(k) = k + \frac{k(k-1)}{2} = \frac{k(k+1)}{2}.
\]
Therefore, the multivariate Gaussian family forms an exponential family of dimension
\[
k + \frac{k(k+1)}{2}.
\]
\end{ex}

\section{Nonuniqueness of Sufficient Statistics and Natural Parameters}%
\label{sec:nssnp}

Exponential-family representations of a statistical model are not unique: one can transform the sufficient statistic and the natural parameter without changing the model. 

\medskip
Let \(T(x)\) be a sufficient statistic and \(\eta(\theta)\) the corresponding natural parameter. For any invertible matrix \(A \in \R^{k \times k}\) and vectors \(\bs{b}, \bs{c} \in \R^k\), define the affine transformations
\[
\tilde{T}(x) = A T(x) + \bs{b}, \qquad 
\tilde{\eta}(\theta) = A^{-T} \eta(\theta) + \bs{c}.
\]

Then the inner product between the transformed parameter and statistic expands as
\begin{align*}
\langle \tilde{\eta}(\theta), \tilde{T}(x) \rangle 
&= \langle A^{-T}\eta(\theta) + \bs{c}, A T(x) + \bs{b} \rangle \\
&\overset{(*)}{=} \langle \eta(\theta), T(x) \rangle 
+ \langle A^{-T}\eta(\theta), \bs{b} \rangle 
+ \langle \bs{c}, A T(x) \rangle 
+ \langle \bs{c}, \bs{b} \rangle
\end{align*}
where in \((*)\) we have use that for any two vectors \(\bs{u }, \bs{v }\) their inner product is define as 
\[
    \left<\bs{u }, \bs{v } \right> = \bs{u }^{T}\bs{v}
\]
so 
\begin{align*}
    \left<A^{-T}\eta(\theta), AT(x) \right>
    &= (A^{-T}\eta(\theta))^{T}AT(x)
    = \eta(\theta)^{T}(A^{-T})^{T}AT(x) 
    = \eta(\theta)^{T}T(x)
    \\
    &= \left<\eta(\theta),T(x) \right>.
\end{align*}

\medskip
This implies that the density can be rewritten as
\[
p_\theta(x)
= \exp\Big\{ \langle \tilde{\eta}(\theta), \tilde{T}(x) \rangle - \tilde{B}(\theta) \Big\} 
  \cdot h(x) \exp\big\{- \langle \bs{c}, A T(x) \rangle\big\},
\]
with the new log-partition function
\[
\tilde{B}(\theta) = B(\theta) + \langle A^{-T} \eta(\theta), \bs{b} \rangle + \langle \bs{c}, \bs{b} \rangle.
\]

\medskip
In other words, an affine transformation of the sufficient statistic \(T(x)\) can be compensated by a corresponding transformation of the natural parameter \(\eta(\theta)\), possibly adjusting the normalizing constant. Linear relations among the coordinates of \(T(x)\) can thus be eliminated by choosing a suitable \(A\) and \(\bs{b}\), leading to a \emph{minimal representation} of the exponential family.  

\medskip
\begin{definition}[Order of an Exponential Family]
The \textbf{order of the exponential family} is the dimension of the minimal sufficient statistic after all linear dependencies among its coordinates have been removed. This dimension is unique for a given family. 
\end{definition}

\medskip
\begin{remark}[Invertibility and singular cases]
    \leavevmode
    \begin{itemize}
        \item 
In univariate or multivariate Gaussian models, assuming \(\Sigma\) positive definite (invertible) ensures the density exists in \(\R^k\). If a variance is zero, the distribution degenerates to a point mass; if a covariance matrix is singular, the random vector takes values in a strict subspace of \(\R^k\). Such singular cases are not part of the exponential-family representation as discussed here.  
\item 
Affine transformations provide a systematic way to expose linear dependencies among coordinates of \(T(x)\). For example, if some linear combination of coordinates is constant across all observations, one can redefine \(T\) via a suitable matrix \(A\) to eliminate the redundant coordinate. This reduction ensures that only the essential dimensions remain, giving the family its unique order.  
    \end{itemize}
\end{remark}

\section{Random Sampling}
\label{sec:rs}

Let \(\mathsf{P}=\{\mathsf{P}_{\theta} : \theta \in \Theta\}\) be a \(k\)--dimensional exponential family with densities
\[
    p_{\theta}(x)
    = \frac{d \mathsf{P}_{\theta}}{d \nu}(x)
    = \exp\{\langle \eta(\theta), T(x) \rangle - B(\theta)\} \, h(x),
    \quad x \in \cc{X}.
\]
Let \(X_1,\dots,X_n\) be i.i.d.\ random variables with distribution \(\mathsf{P}_{\theta}\).

The joint distribution of the random vector \(\bs{X}=(X_1,\dots,X_n)\) then has density
\[
    p_{\theta}(\bs{x})
    = \prod\limits_{i=1}^{n} p_{\theta}(x_i)
    = \exp\left\{
        \Big\langle \eta(\theta), \sum\limits_{i=1}^{n} T(x_i) \Big\rangle
        - n B(\theta)
      \right\}
      \cdot
      \prod\limits_{i=1}^{n} h(x_i),
    \quad \bs{x} \in \cc{X}^n.
\]

Thus, the family of joint distributions of \(\bs{X}\) is again a \(k\)--dimensional exponential family.
The natural parameter \(\eta(\theta)\) remains unchanged, while the sufficient statistic becomes
\[
    \sum\limits_{i=1}^{n} T(X_i).
\]
In particular, independent sampling preserves the exponential family structure.

\begin{note}
The dimension of the sufficient statistic
does not depend on the sample size \(n\).\footnote{This behavior is highly atypical outside the class
of exponential families; see Hipp (1974). Recall, however, the \(\mathrm{Uniform}(0,\theta)\) model
as a counterexample without regularity.}
\end{note}

This property has important consequences. Even for very large samples, all relevant statistical
information is summarized by a single vector in \(\mathbb{R}^k\). As a result, statistical analysis,
geometric intuition, and asymptotic arguments for exponential families can be carried out in a fixed
finite-dimensional space, regardless of the number of observations.

By contrast, for general statistical models, sufficient statistics may grow in dimension with the
sample size (for example, order statistics). Under mild regularity conditions, the invariance of
the dimension of sufficient statistics under i.i.d.\ sampling essentially characterizes exponential
families.

\section{Minimal Sufficiency}
\label{sec:Minimal Sufficiency}

For exponential families, minimal sufficiency can be characterized in a particularly transparent
way. The key reason is that likelihood ratios in exponential families reduce to differences of
linear forms in the sufficient statistic. As a result, the general criterion for minimal sufficiency
based on proportional likelihoods admits a simple geometric interpretation in terms of the natural
parameter space; see Corollary~6.16 in Lehmann and Casella (1998).

\begin{corollary}
    Let \(\mathsf{P}=\{\mathsf{P}_{\theta} : \theta \in \Theta\}\) be a \(k\)--dimensional exponential
    family with natural parameter mapping \(\eta : \Theta \to \R^k\).
    If the set
    \[
        \eta(\Theta) = \{\eta(\theta) : \theta \in \Theta\} \subseteq \R^k
    \]
    contains \(k+1\) points \(\eta^{(0)},\dots,\eta^{(k)}\) such that the vectors
    \[
        \eta^{(j)} - \eta^{(0)}, \quad j=1,\dots,k,
    \]
    are linearly independent, then the sufficient statistic \(T\) is minimal sufficient.
\end{corollary}

This result follows directly from the characterization of minimal sufficiency in terms of
proportional likelihoods. In an exponential family, the ratio of two likelihoods corresponding
to parameter values \(\theta\) and \(\theta'\) can be written as
\[
    \exp\Big\{
        \big\langle \eta(\theta) - \eta(\theta'),\, T(x) \big\rangle
        - \big(B(\theta) - B(\theta')\big)
    \Big\},
\]
so equality of likelihood ratios across all parameter values is determined entirely by linear
constraints on \(T(x)\). If the set of possible difference vectors
\(\eta(\theta) - \eta(\theta')\) spans \(\R^k\), then no nontrivial reduction of \(T\) can preserve
these ratios, and \(T\) is therefore minimal sufficient.

An equivalent geometric formulation of the condition in the corollary is the following.
The sufficient statistic \(T\) is minimal sufficient if and only if the affine hull of
\(\eta(\Theta)\) has dimension \(k\).
The affine hull of a set \(A \subseteq \R^k\) is defined as
\[
    \left\{
        \sum\limits_{j=0}^{k} \alpha_j \eta^{(j)} :
        \eta^{(j)} \in A,\;
        \alpha_j \in \R,\;
        \sum\limits_{j=0}^{k} \alpha_j = 1
    \right\}.
\]

In many standard examples, this condition is easily verified.

\paragraph{Univariate Gaussian family.}
Consider the normal distribution with unknown mean \(\mu \in \R\) and variance
\(\sigma^2 > 0\).
The parameter space \(\Theta = \R \times (0,\infty)\) is the open upper half--plane.
The corresponding natural parameter space \(\eta(\Theta)\) is also an open subset of
\(\R^2\), hence its affine hull has dimension \(2\).
Therefore, the canonical sufficient statistic is minimal sufficient.

\paragraph{Multivariate Gaussian family.}
For the multivariate normal distribution with mean vector \(\mu \in \R^k\) and covariance
matrix \(\Sigma\) positive definite, the natural parameter space is
\(\R^k \times \mathcal{S}_{++}^k\), where \(\mathcal{S}_{++}^k\) denotes the cone of positive
definite matrices.
This set is open and full--dimensional in the corresponding Euclidean space, implying minimal
sufficiency of the usual sufficient statistic.

\paragraph{Binomial family and logistic regression.}
For the binomial model, the original parameter \(\theta\) lies in the open unit interval
\((0,1)\).
The natural parameter is given by the log--odds transformation
\[
    \eta(\theta) = \log\frac{\theta}{1-\theta},
\]
which maps \((0,1)\) onto all of \(\R\).
Consequently, \(\eta(\Theta) = \R\), whose affine hull has dimension \(1\).
This explains both the minimal sufficiency of the binomial sufficient statistic and the
central role of the log--odds parameterization in logistic regression, where linear predictors
naturally take values in \(\R\).

In summary, for exponential families, minimal sufficiency is governed entirely by the geometry
of the natural parameter space.
If \(\eta(\Theta)\) is full--dimensional (or equivalently, has an affine hull of dimension \(k\)),
then the canonical sufficient statistic is automatically minimal sufficient.

%For instance, in the Gaussian location--scale family, the natural parameter space is an open
%subset of \(\R^2\); in the binomial model, the natural parameter (log--odds) ranges over all of
%\(\R\).
%In such cases, the affine hull of \(\eta(\Theta)\) is full--dimensional, and the canonical
%sufficient statistic is automatically minimal sufficient.
%
%Thus, under mild regularity conditions, minimal sufficiency in exponential families is governed
%entirely by the geometry of the natural parameter space.

\section{Canonical Form}
\label{sec:Canonical Form}

By re--parametrizing an exponential family in terms of its natural parameters, the model can be
written in \textbf{canonical form}. Specifically, a \(k\)--dimensional exponential family admits
a representation
\[
    p_{\eta}(x)
    = \exp\{\langle \eta, T(x) \rangle - A(\eta)\} \, h(x),
    \quad x \in \cc{X},
\]
where \(\eta \in \R^k\) is the natural parameter and \(T(x) \in \R^k\) is the sufficient statistic.
We write \(\mathsf{P}_{\eta}\) for the distribution with density \(p_{\eta}\).

Passing to the canonical form amounts to abandoning the original parameterization
(e.g.\ means, variances, probabilities, rates) in favor of the natural parameters.
For instance, instead of indexing a Gaussian distribution by its mean and variance,
or a binomial distribution by a success probability, we index these distributions by
inverse variances, mean--variance ratios, or log--odds.
This is purely a re--parametrization: the underlying distributions do not change, only the
coordinates used to describe them.

\subsection{Natural Parameter Space}

The density \(p_{\eta}(x)\) is well defined only for those values of \(\eta\) for which the
normalizing constant is finite.
This leads to the definition of the \textbf{natural parameter space}
\[
    \cc{H}
    =
    \left\{
        \eta \in \R^k :
        \int_{\cc{X}} \exp\{\langle \eta, T(x) \rangle\} h(x) \, d\nu(x) < \infty
    \right\}.
\]
For every \(\eta \in \cc{H}\), a valid probability density is obtained by setting
\[
    A(\eta)
    =
    \log
    \int_{\cc{X}}
        \exp\{\langle \eta, T(x) \rangle\} h(x) \, d\nu(x).
\]
The function \(A(\eta)\) is called the \textbf{cumulant generating function} or
\textbf{log--Laplace transform} of the family.

The role of \(\cc{H}\) is fundamental.
Any nonnegative function with finite integral can be normalized to obtain a probability
density, but if the integral diverges, no normalization is possible.
Thus, \(\cc{H}\) is the largest set of natural parameters for which the exponential
representation yields a well-defined statistical model.

\subsection*{Examples}

\paragraph{Gaussian family.}
For the univariate Gaussian distribution, a natural sufficient statistic is
\(T(x) = (x, x^2)\).
Only certain choices of \(\eta\) lead to integrable densities:
the coefficient multiplying \(x^2\) must be negative.
This restriction corresponds exactly to the requirement that the variance be positive.
If the quadratic term had the wrong sign, the exponential would not be integrable over
\(\R\), and no probability density could be defined.

\paragraph{Binomial family and logistic regression.}
For the binomial model with success probability \(\theta \in (0,1)\), the natural parameter is
the log--odds
\[
    \eta = \log\frac{\theta}{1-\theta},
\]
which ranges over all of \(\R\).
Thus, the natural parameter space is unconstrained, even though the original parameter
\(\theta\) lies in a bounded interval.
This re--parametrization explains why logistic regression works with linear predictors:
the natural parameter lives on the entire real line, making linear combinations natural
and unconstrained.

\paragraph{Multinomial and symmetry considerations.}
In multinomial models, it is sometimes convenient to retain a redundant representation of
the sufficient statistic in order to preserve symmetry.
While redundant coordinates may later be removed to obtain a minimal representation,
keeping them can simplify calculations and notation without affecting the underlying
statistical structure.

\paragraph{Exponential random graph models.}
The exponential family framework can also be used constructively to define models.
For example, when modeling random graphs, the observation \(X\) is a graph, and one may
choose a sufficient statistic \(T(X)\) that counts features such as the number of edges
and the number of triangles.
Given parameters \(\eta_1,\eta_2\), the probability of observing a particular graph is
proportional to
\[
    \exp\{\eta_1 T_1(X) + \eta_2 T_2(X)\}.
\]
Provided the graph space is finite, the normalizing constant \(A(\eta)\) always exists.
Varying \(\eta\) yields a flexible family of distributions that can encode complex
network structure beyond independent edge formation.

\subsection*{Convexity and Likelihood Optimization}

A key advantage of the canonical form emerges when studying likelihood functions.
In general statistical models, likelihoods may be highly non--convex and multimodal.
In contrast, for exponential families in canonical form, the log--likelihood is a concave
function of \(\eta\) over the natural parameter space \(\cc{H}\).

This explains a common empirical observation:
for many standard models (binomial, multinomial, Poisson, Gaussian, gamma),
maximum likelihood estimators are unique and easy to compute.
Although this fact may appear model--specific when working in the original parameters,
it has a unified explanation in the canonical parameterization.

\subsection*{Moment and Cumulant Generating Functions}

Finally, note that without the logarithm, the integral defining \(A(\eta)\) corresponds to
a moment generating function.
For a random variable \(X\), the moment generating function is
\(\mathsf{E}[e^{tX}]\);
for a random vector, it becomes
\(\mathsf{E}[\exp(\langle t, X \rangle)]\).
Taking the logarithm yields the cumulant generating function, whose derivatives generate
moments and cumulants of the sufficient statistic.

Thus, the canonical form not only simplifies likelihood theory and optimization, but also
provides a natural bridge between exponential families, convex analysis, and moment theory.

\subsection{Convexity Properties}
\label{sub:Convexity Properties}

One of the most powerful structural features of exponential families is their built--in
convexity. This convexity governs both the geometry of the parameter space and the shape of
likelihood functions, and it explains why maximum likelihood estimation is particularly
well behaved in these models.

Recall that in canonical form an exponential family has densities
\[
    p_{\eta}(x)
    =
    \exp\{\langle \eta, T(x) \rangle - A(\eta)\} h(x),
    \quad \eta \in \cc{H},
\]
where the natural parameter space \(\cc{H}\) consists precisely of those \(\eta\) for which
\(A(\eta) < \infty\).

\begin{theo}[Convexity Properties]
\label{theo:Convexity}
\leavevmode
\begin{enumerate}[label=(\roman*)]
    \item The natural parameter space \(\cc{H}\) is a convex subset of \(\R^k\).
    \item The cumulant generating function \(A(\eta)\) is convex on \(\cc{H}\).
        \label{ite:2}
    It is strictly convex if
    \(\mathsf{P}_{\eta} \neq \mathsf{P}_{\eta'}\) for all
    \(\eta,\eta' \in \cc{H}\) with \(\eta \neq \eta'\).
    \item For fixed data \(x\), the log--likelihood function
        \label{ite:3}
    \[
        \ell_x(\eta) = \log p_{\eta}(x)
    \]
    is concave on \(\cc{H}\).
    It is strictly concave if
    \(\mathsf{P}_{\eta} \neq \mathsf{P}_{\eta'}\) for all
    \(\eta,\eta' \in \cc{H}\) with \(\eta \neq \eta'\).
\end{enumerate}
\end{theo}

\myparagraph{Comments:}
Convexity is the mathematical reason why maximum likelihood estimation in exponential
families rarely exhibits pathologies such as multiple local optima.
In contrast to general statistical models, likelihood optimization here is a convex
optimization problem.

Recall the basic facts:
\begin{itemize}
    \item Every local maximum of a concave function is a global maximum.
    \item A strictly concave function has at most one maximizer.
\end{itemize}
Consequently, in the \emph{full} exponential family indexed by all
\(\eta \in \cc{H}\), the maximum likelihood estimator (if it exists) is unique.

Existence, however, is not guaranteed: a strictly concave function may fail to attain its
supremum (for example, the logarithm on \((0,\infty)\)).
Nevertheless, uniqueness follows automatically once existence is established.

If one restricts \(\eta\) to an arbitrary subset of \(\cc{H}\), convexity and concavity
may be destroyed.
An important exception occurs when the restriction is a linear subspace of \(\cc{H}\):
in that case, the model remains an exponential family of lower dimension, and the theorem
continues to apply.

\begin{proof}[]
We begin with convexity of the natural parameter space and of the function \(A\).
Let \(\eta, \eta' \in \cc{H}\) and \(\alpha \in [0,1]\).
By definition,
\begin{align*}
    \exp\{A(\alpha \eta + (1-\alpha)\eta')\}
    &=
    \int
        \exp\{\langle \alpha \eta + (1-\alpha)\eta', T(x) \rangle\}
        h(x)\, d\nu(x) \\
    &=
    \int
        \left( \exp\{\langle \eta, T(x) \rangle\} \right)^{\alpha}
        \left( \exp\{\langle \eta', T(x) \rangle\} \right)^{1-\alpha}
        h(x)\, d\nu(x).
\end{align*}
This expression is exactly of the form to which Hölder's inequality applies.
Using Hölder's inequality yields
\begin{align*}
    \exp\{A(\alpha \eta + (1-\alpha)\eta')\}
    &\le
    \left(
        \int \exp\{\langle \eta, T(x) \rangle\} h(x)\, d\nu(x)
    \right)^{\alpha}
    \left(
        \int \exp\{\langle \eta', T(x) \rangle\} h(x)\, d\nu(x)
    \right)^{1-\alpha} \\
    &=
    \exp\{\alpha A(\eta) + (1-\alpha) A(\eta')\}.
\end{align*}
Taking logarithms gives
\begin{equation}
    \label{eq:convex}
    A(\alpha \eta + (1-\alpha)\eta')
    \le
    \alpha A(\eta) + (1-\alpha) A(\eta').
\end{equation}

This inequality has two immediate consequences.
First, if \(A(\eta)\) and \(A(\eta')\) are finite, then
\(A(\alpha \eta + (1-\alpha)\eta')\) is finite as well.
Hence, \(\cc{H}\) is convex.
Second, inequality \eqref{eq:convex} is precisely the definition of convexity of the function
\(A\).

Equality in Hölder's inequality occurs if and only if the two functions inside the integral
are proportional.
Here, this would require the existence of a constant \(c > 0\) such that
\[
    \exp\{\langle \eta, T(x) \rangle\}
    =
    c \cdot \exp\{\langle \eta', T(x) \rangle\}
    \quad \text{for all } x.
\]
After normalization, this implies
\(\mathsf{P}_{\eta} = \mathsf{P}_{\eta'}\).
Therefore, if distinct natural parameters always correspond to distinct distributions,
the inequality in \eqref{eq:convex} is strict for \(\alpha \in (0,1)\), and \(A\) is strictly
convex, proving \ref{ite:2}.

Finally, for fixed data \(x\),
\[
    \ell_x(\eta)
    =
    \log p_{\eta}(x)
    =
    \langle \eta, T(x) \rangle - A(\eta) + \log h(x).
\]
The term \(\langle \eta, T(x) \rangle\) is linear in \(\eta\), while \(-A(\eta)\) is concave
because \(A\) is convex.
Adding a linear function does not affect concavity.
Hence, \(\ell_x(\eta)\) is concave on \(\cc{H}\), and strictly concave whenever \(A\) is
strictly convex, proving \ref{ite:3}.
\end{proof}

The convexity of the natural parameter space and the cumulant generating function provides a
unified explanation for the uniqueness of maximum likelihood estimators in exponential
families such as the Gaussian, Poisson, binomial, and multinomial models.
What appears as a collection of model--specific calculations is, in fact, a single geometric
phenomenon.


\begin{ex}[Gaussian model and convexity of the cumulant function]
Let \(X_1,\dots,X_n\) be i.i.d.\ random variables with
\[
X_i \sim \mathcal{N}(\mu,\sigma^2),
\]
where both the mean \(\mu\) and the variance \(\sigma^2\) are unknown.  
We represent this model as a full exponential family by introducing the natural
parameters
\[
\eta_1 = \frac{\mu}{\sigma^2},
\qquad
\eta_2 = \frac{1}{\sigma^2}.
\]

With this parametrization, the log-likelihood function can be written in the
canonical exponential-family form
\[
\ell_x(\eta)
= \eta_1 \sum_{i=1}^n x_i
  + \eta_2 \left(-\frac{1}{2}\sum_{i=1}^n x_i^2 \right)
  - A(\eta),
\]
where the cumulant generating function (log-partition function) is given by
\[
A(\eta)
= \frac{n}{2}\frac{\mu^2}{\sigma^2}
  + \frac{n}{2}\log(2\pi\sigma^2)
= \frac{n}{2}\left(
    \frac{\eta_1^2}{\eta_2}
    - \log \eta_2
    + \log(2\pi)
  \right).
\]

The gradient of \(A\) coincides with the vector of expectations of the sufficient
statistics. A direct calculation yields
\[
\nabla A(\eta)
=
\begin{pmatrix}
\displaystyle n\,\frac{\eta_1}{\eta_2} \\[0.6em]
\displaystyle -\frac{n}{2}\left(
  \frac{\eta_1^2}{\eta_2^2} + \frac{1}{\eta_2}
\right)
\end{pmatrix}
=
\begin{pmatrix}
n\mu \\[0.4em]
-\frac{n}{2}\bigl(\mu^2+\sigma^2\bigr)
\end{pmatrix}
=
\begin{pmatrix}
\mathsf{E}_\eta\!\left[\sum\limits_{i=1}^n X_i\right] \\[0.6em]
\mathsf{E}_\eta\!\left[-\frac{1}{2}\sum\limits_{i=1}^n X_i^2\right]
\end{pmatrix}.
\]
Thus, as predicted by the general theory of exponential families, the gradient of
the cumulant function returns the expectations of the sufficient statistics.

To verify the convexity of \(A\) directly, we compute its Hessian matrix:
\[
\nabla^2 A(\eta)
=
\begin{pmatrix}
\displaystyle \frac{n}{\eta_2}
&
\displaystyle -\,\frac{n\eta_1}{\eta_2^2}
\\[0.8em]
\displaystyle -\,\frac{n\eta_1}{\eta_2^2}
&
\displaystyle n\!\left(
  \frac{\eta_1^2}{\eta_2^3}
  + \frac{1}{2\eta_2^2}
\right)
\end{pmatrix}.
\]
Since \(\eta_2>0\) (it is the inverse variance), the upper-left entry is positive.
Moreover, the determinant is
\[
\det\bigl(\nabla^2 A(\eta)\bigr)
= \frac{n^2}{2\eta_2^3} > 0.
\]
By Sylvester’s criterion, the Hessian is positive definite for all admissible
\(\eta\), and hence \(A\) is strictly convex.

This explicit calculation illustrates the general convexity result for cumulant
functions in exponential families, which in theory follows from Hölder’s
inequality but in this Gaussian case can be verified directly by elementary
calculus.
\end{ex}

\subsection{Analytic Properties}
\label{sub:ap}

\begin{theo}[Analyticity of exponential-family integrals]
\label{theo:ap}
Let \(\eta\) be a point in the interior of the natural parameter space \(\cc{H}\).
Let \(f : \cc{X} \to \R\) be a measurable function that is integrable with respect
to \(\mathsf{P}_{\eta}\).

Then the mapping
\[
\eta \longmapsto
\int_{}^{} f(x)\,
\exp\!\bigl\{\langle \eta, T(x) \rangle\bigr\}
\, h(x)\, d\nu(x)
\]
admits continuous partial derivatives of all orders. Moreover, all derivatives
can be computed by differentiating under the integral sign. In particular, the
function is analytic on the interior of \(\cc{H}\).
\end{theo}

\begin{remark}
The expression above represents (up to normalization) the expectation of the
statistic \(f(X)\) when \(X\) is drawn from an exponential family distribution
with natural parameter \(\eta\). The theorem therefore shows that expectations
with respect to exponential family models depend smoothly on the natural
parameters, as long as \(\eta\) lies in the interior of the parameter space.
\end{remark}

\begin{remark}[Idea of the proof]
The result follows from an application of Lebesgue’s dominated convergence
theorem to difference quotients arising in the definition of partial
derivatives. Since derivatives are limits of difference quotients, the key step
is to justify the interchange of differentiation and integration. The
exponential structure of the density allows the difference quotients to be
dominated by an integrable function, ensuring the validity of this
interchange. A detailed proof can be found in Lehmann and Romano (2005,
Theorem~2.7.1).
\end{remark}

\subsection{Cumulant Generating Transform}
\label{sub:cgt}

A central object in the theory of exponential families is the cumulant generating
function \(A(\eta)\). Beyond its role as a normalizing constant, \(A\) encodes the
first- and second-order moments of the sufficient statistic through
differentiation.

\begin{corollary}[Derivatives of the cumulant function]
\label{cor:cgt}
Let \(\{ \mathsf{P}_\eta : \eta \in \cc{H} \}\) be an exponential family in canonical
form with sufficient statistic
\(
T = (T_1,\dots,T_d)
\)
and cumulant generating function \(A\).
For all \(\eta\) in the interior of \(\cc{H}\),
\[
\frac{\partial}{\partial \eta_j} A(\eta)
= \mathsf{E}_\eta\!\left[T_j(X)\right],
\qquad
\frac{\partial^2}{\partial \eta_j \partial \eta_k} A(\eta)
= \mathsf{Cov}_\eta\!\left[T_j(X),T_k(X)\right].
\]
Equivalently,
\[
\nabla A(\eta) = \mathsf{E}_\eta[T(X)],
\qquad
\nabla^2 A(\eta) = \mathsf{Cov}_\eta[T(X)].
\]
\end{corollary}

\begin{proof}
By definition of the exponential family in canonical form, the density integrates
to one:
\[
\int_{\cc{X}}
\exp\!\left\{
\langle \eta, T(x) \rangle - A(\eta)
\right\}
h(x)\, d\nu(x)
= 1,
\qquad \forall\, \eta \in \cc{H}.
\]
Although the right-hand side is constant, the left-hand side is a function of
\(\eta\). Since \(\eta\) lies in the interior of \(\cc{H}\), we may differentiate
under the integral sign.

Taking the partial derivative with respect to \(\eta_j\) yields
\[
\int_{\cc{X}}
\exp\!\left\{
\langle \eta, T(x) \rangle - A(\eta)
\right\}
\left(
T_j(x) - \frac{\partial}{\partial \eta_j} A(\eta)
\right)
h(x)\, d\nu(x)
= 0.
\]
Recognizing the integrand as an expectation under \(\mathsf{P}_\eta\) proves the
first identity.



Differentiating once more with respect to \(\eta_k\) 
and differentiating under the
integral sign, yields
\begin{align*}
0
&=
\int_{\cc{X}}
\exp\!\left\{
\langle \eta, T(x) \rangle - A(\eta)
\right\}
\Bigl(
T_k(x) - \frac{\partial}{\partial \eta_k} A(\eta)
\Bigr)
\Bigl(
T_j(x) - \frac{\partial}{\partial \eta_j} A(\eta)
\Bigr)
h(x)\, d\nu(x)
\\
&\qquad
-
\int_{\cc{X}}
\exp\!\left\{
\langle \eta, T(x) \rangle - A(\eta)
\right\}
\frac{\partial^2}{\partial \eta_j \partial \eta_k} A(\eta)
\, h(x)\, d\nu(x).
\end{align*}
Since the density integrates to one, the second integral equals
\(\frac{\partial^2}{\partial \eta_j \partial \eta_k} A(\eta)\), and we obtain
\[
\int_{\cc{X}}
\exp\!\left\{
\langle \eta, T(x) \rangle - A(\eta)
\right\}
\Bigl(
T_j(x) - \mathsf{E}_\eta[T_j(X)]
\Bigr)
\Bigl(
T_k(x) - \mathsf{E}_\eta[T_k(X)]
\Bigr)
h(x)\, d\nu(x)
=
\frac{\partial^2}{\partial \eta_j \partial \eta_k} A(\eta).
\]
Recognizing the left-hand side as a covariance under \(\mathsf{P}_\eta\), we
conclude that
\[
\mathsf{Cov}_\eta\!\left[T_j(X),T_k(X)\right]
=
\frac{\partial^2}{\partial \eta_j \partial \eta_k} A(\eta).
\]

\end{proof}

\begin{remark}[Interpretation]
The function \(A\) is called the cumulant generating function because its
derivatives generate cumulants of the sufficient statistic. In particular, the
gradient of \(A\) yields the mean vector, while the Hessian yields the covariance
matrix. This provides a direct and efficient way to compute moments in exponential
families by differentiation rather than integration.
\end{remark}

\subsection{Random Vectors}
\label{sub:rv}

Let
\[
X = \begin{pmatrix}
X_1 \\ \vdots \\ X_d
\end{pmatrix}
\]
be a random vector, viewed equivalently as a vector of \(d\) real-valued random
variables or as a random element of \(\R^d\) equipped with its Borel
\(\sigma\)-algebra.

\begin{definition}[Expectation of a random vector]
The expectation (mean vector) of \(X\) is defined componentwise by
\[
\mathsf{E}[X]
=
\begin{pmatrix}
\mathsf{E}[X_1] \\ \vdots \\ \mathsf{E}[X_d]
\end{pmatrix},
\]
whenever all component expectations exist.
\end{definition}

\begin{definition}[Variance matrix]
The variance matrix of \(X\) is the \(d \times d\) matrix
\[
\mathsf{Var}[X]
=
\bigl(\mathsf{Cov}(X_i,X_j)\bigr)_{1 \le i,j \le d},
\]
whose diagonal entries are the variances of the coordinates of \(X\) and whose
off-diagonal entries are the pairwise variances.
\end{definition}

\begin{remark}
Expectations of vector- or matrix-valued random quantities are always understood
entrywise. In particular, whenever \(\mathsf{E}[X]\) or \(\mathsf{Var}[X]\) is
written, all relevant moments are assumed to exist.
\end{remark}

\begin{remark}[Positive semidefiniteness]
The covariance matrix \(\mathsf{Var}[X]\) is always positive semidefinite. 
\end{remark}

In the context of exponential families, the sufficient statistic \(T(X)\) is a
random vector. Applying the results of
Section~\ref{sub:cgt}, the gradient and Hessian of the cumulant generating function
\(A\) admit a direct probabilistic interpretation:
\[
\nabla A(\eta)
=
\mathsf{E}_\eta[T(X)],
\qquad
D_2 A(\eta)
=
\mathsf{Var}_\eta[T(X)].
\]
Thus, the first derivative of \(A\) yields the mean vector of the sufficient
statistic, while the second derivative yields its variance matrix.

\section{Mean Parametrization}
\label{sec:mp}

Let \(\{\mathsf{P}_\eta : \eta \in \cc{H}'\}\) be an exponential family in canonical
form, where the natural parameter space \(\cc{H}' \subset \R^d\) is open and
convex. The associated cumulant generating function \(A : \cc{H}' \to \R\) is
strictly convex on \(\cc{H}'\).

As a consequence, the gradient mapping
\[
\eta \longmapsto \nabla A(\eta)
=
\mathsf{E}_\eta[T(X)]
\]
is injective. Hence, each value of the natural parameter \(\eta\) corresponds to
a unique value of the expectation of the sufficient statistic. This establishes a
one--to--one correspondence between the natural parameter \(\eta\) and the
\emph{mean parameter}
\[
\theta \;\coloneqq\; \mathsf{E}_\eta[T(X)].
\]

This alternative parametrization of the exponential family is called the
\emph{mean parametrization}. While the natural parametrization is particularly
well-suited for studying convexity and optimization properties, the mean
parametrization provides a direct probabilistic interpretation in terms of
expected sufficient statistics.

\begin{ex}[Gaussian model]
Consider the Gaussian exponential family with unknown mean \(\mu\) and variance
\(\sigma^2\), and natural parameters
\[
\eta
=
\begin{pmatrix}
\eta_1 \\ \eta_2
\end{pmatrix}
=
\begin{pmatrix}
\mu / \sigma^2 \\ 1 / \sigma^2
\end{pmatrix}.
\]
The corresponding mean parameters are given by
\[
\nabla A(\eta)
=
\mathsf{E}_\eta[T(X)]
=
\begin{pmatrix}
\mu \\[0.4em]
-\frac{1}{2}(\mu^2 + \sigma^2)
\end{pmatrix}
=
\begin{pmatrix}
\displaystyle \frac{\eta_1}{\eta_2} \\[0.8em]
\displaystyle -\frac{1}{2}\left(
\frac{\eta_1^2}{\eta_2^2} + \frac{1}{\eta_2}
\right)
\end{pmatrix}.
\]
This mapping is one--to--one, illustrating explicitly the equivalence between the
natural and mean parametrizations in this model.
\end{ex}

\begin{remark}
The injectivity of the gradient map follows from strict convexity of \(A\): the
gradient of a strictly convex function is injective on any convex domain.
Consequently, distributions in an exponential family may be uniquely identified
either by their natural parameters or by the mean values of their sufficient
statistics.
\end{remark}

\section{A Terminological Remark on Exponential Families}

Before turning to estimation, we make a brief remark on terminology. An
\emph{exponential family} is not a single probability distribution but a
\emph{statistical model}, that is, a collection of probability distributions
indexed by a parameter. Typical examples include the Gaussian distributions with
varying mean and variance, the Poisson distributions with varying rate parameter,
or the binomial distributions with varying success probability.

In the literature—particularly in machine learning—it is common to encounter
phrases such as “the Poisson distribution is in the exponential family.” Such
formulations should be interpreted with care. More precisely, one should say
that a \emph{family} of Poisson distributions, indexed by their rate parameter,
\emph{forms} or \emph{constitutes} an exponential family when written in
canonical form.

Throughout these notes, we will therefore avoid referring to individual
distributions as belonging to “the” exponential family. Instead, we emphasize
that exponential families are collections of distributions that admit a common
representation in exponential-family form.

We now proceed to the problem of statistical estimation for exponential family
models.
